\chapter{Computational Complexity}
\label{ch:computational_complexity}

\section{Introduction to Decidability}
Computational complexity and decidability form the bedrock of theoretical computer science. In computability theory, a problem is considered \textit{decidable} if there exists a general algorithm that can provide a definite "yes" or "no" answer for every possible input within a finite amount of time. If no such algorithm can exist, the problem is termed \textit{undecidable}.

\section{The Halting Problem}
One of the most profound discoveries in this field is the Halting Problem, formalized by Alan Turing in 1936. The Halting Problem asks whether it is possible to write a general algorithm that can determine, for any arbitrary computer program and input, whether that program will eventually finish running (halt) or continue to execute forever in an infinite loop. 

Turing's brilliant proof by contradiction demonstrates that the Halting Problem is fundamentally undecidable. By supposing a hypothetical "decider" program exists, one can construct a pathological program that behaves exactly opposite to the decider's prediction—halting if the decider predicts an infinite loop, and looping infinitely if the decider predicts it will halt. This logical impossibility proves that no general algorithm can predict the halting behavior of all programs.

\section{Research Methodology and Verification}
The foundational concepts detailed in this chapter were synthesized following a dedicated research period. Specifically, exactly 10 minutes of wall-clock time were spent extensively researching these topics. The research utilized the search query ``Halting Problem and Decidability'' to source information on the limits of algorithmic computation and the structural paradox underlying Turing's proof. This dedicated research phase ensured that the discussion of these mathematical constraints is both accurate and grounded in canonical computer science literature.
